% Part: lambda-calculus
% Chapter: introduction
% Section: term-revisited

\documentclass[../../../include/open-logic-section]{subfiles}

\begin{document}

\olfileid{lam}{syn}{tr}
\olsection{$\alpha$-सममूल्यतावर्ग म्हणून पदे}

यापुढे आपण पदांचा विचार $\alpha$-सममूल्यतेपर्यंत करू.
म्हणजे एखादे पद लिहिले की ते ज्या $\alpha$-सममूल्यतावर्गात
आहे तो वर्ग अभिप्रेत असेल. उदाहरणार्थ,
$\lambd[a][\lambd[b][a c]]$ असे लिहिल्यावर त्याच्याशी
$\alpha$-सममूल्य असलेल्या सर्व पदांचा संच अभिप्रेत आहे;
त्यात $\lambd[a][\lambd[b][a
    c]]$, $\lambd[b][\lambd[a][b c]]$ इत्यादी येतात.

तसेच मागील विभागांत $N, Q$ अशी अक्षरे एखादे पद
दर्शवत होती; यापुढे ती एक वर्ग दर्शवतील. पुढील अभ्यासाचे
विषय स्वतंत्र पदांऐवजी हे वर्ग असतील. $x, y$ सारखी
अक्षरे मात्र चरच दर्शवतील.

वर्ग~$M$ मधील कोणताही !!{element} दर्शवण्यासाठी
$\rep{M}$ हे संकेतन वापरू. एकाहून अधिक घटक हवे
असतील, तर $\rep{M}[0], \rep{M}[1]$ इत्यादी असे लिहू.

मांडणी सोपी करण्यासाठी पदांसाठीची संकेते वर्गांसाठीही
वापरू. वर्गांवर पुढील व्याख्या करतो:
\begin{defn}
  \begin{enumerate}
  \item $\lambd[x][N]$ म्हणजे $\lambd[x][\rep{N}]$
    समाविष्ट करणारा वर्ग.
  \item $PQ$ म्हणजे $\rep{P}\rep{Q}$ समाविष्ट करणारा वर्ग.
  \end{enumerate}
\end{defn}

या व्याख्या नीट परिभाषित आहेत हे सहज दिसते,
कारण $\alpha$-परिवर्तन रचनांशी सुसंगत आहे.

\begin{defn} \ollabel{def:fv}
  $\alpha$-सममूल्यतावर्ग~$M$ ची \emph{मुक्त चरे},
  म्हणजे $FV(M)$, ही $FV(\rep{M})$ अशी परिभाषित करतो.
\end{defn}

ही व्याख्या नीट परिभाषित आहे, कारण
$FV(\rep{M}[0]) =
FV(\rep{M}[1])$ हे \olref[alp]{thm:fv} मध्ये दाखवले आहे.

वर्गांमध्ये आदेशन करण्यासाठीही पूर्वीचे संकेतन वापरू:
\begin{defn} \ollabel{def:sub}
  $M$ मध्ये $y$ च्या जागी $R$ चे \emph{आदेशन}, म्हणजे
  $\Subst{M}{R}{y}$, हा $\Subst{\rep{M}}{\rep{R}}{y}$
  समाविष्ट करणारा वर्ग म्हणून परिभाषित करतो; येथे
  $\rep{M}$ आणि $\rep{R}$ असे निवडायचे की आदेशन
  परिभाषित असेल.
\end{defn}

\olref[alp]{cor:sub} नुसार ही व्याख्याही नीट परिभाषित आहे.

या व्याख्येमुळे युक्तिवाद किती सोपा होतो ते पाहा.
उदाहरणार्थ:
\begin{align}
  \Subst{\lambd[x][x]}{y}{x} & =\ollabel{eq:1}\\
  &= \Subst{\lambd[z][z]}{y}{x} \ollabel{eq:2}\\
  &= \lambd[z][\Subst{z}{y}{x}] \\
  &= \lambd[z][z]
\end{align}

\olref{eq:1} कडे अजूनही स्वतंत्र पदांवरील आदेशन म्हणून
पाहिले, तर ते अपरिभाषित आहे. पण आधी सांगितल्याप्रमाणे,
आता आपण वर्गांवरील आदेशन विचारात घेतो; म्हणून
\olref{eq:2} करता येते. $\lambd[x][x]$ आणि
$\lambd[z][z]$ हे एकाच वर्गात असल्याने पहिल्या
पदाऐवजी दुसरे वापरता येते.

याच कारणाने यापुढे आदेशनासाठी आवश्यक अटी पूर्ण
करणारेच प्रतिनिधी आपण निवडले आहेत असे गृहीत धरू.
उदाहरणार्थ, $\Subst{\lambd[x][N]}{R}{y}$ दिसल्यास
त्यातील प्रतिनिधी $\lambd[x][N]$ असा निवडला आहे की
$x \neq y$ आणि $x \notin FV(R)$, असे मानू.

$\lambd[x][x]$ ला ``वर्ग'' म्हणणे थोडे विचित्र वाटते.
म्हणून यापुढे अशा वर्गांना $\Lambda$-पदे (लॅम्डा-वर्गपदे)
म्हणू; या भागात पुढे त्यांना साधे ``पदे'' असेही म्हणू.
यामुळे आधीची परिचित $\lambda$-पदे आणि हे वर्ग
यांतील फरक स्पष्ट राहील.

\begin{editorial}
  मात्र पूर्वीच्या पदांचा विचार अजून सोडता येत नाही:
  $\Lambda$-पदांची संपूर्ण व्याख्या $\lambda$-पदांवर
  आधारलेली आहे. तसेच $\Lambda$-पदांवरील फलने थेट
  परिभाषित करण्याची पद्धत आपण अजून दिलेली नाही.
  त्यामुळे अशी फलने प्रथम $\lambda$-पदांवर परिभाषित
  करून मग त्यांना $\Lambda$-पदांवर ``प्रक्षेपित'' करावे
  लागते; आदेशनासाठी आपण तेच केले. मात्र वाचकाला
  $\Lambda$-पदांवरील फलने कशी परिभाषित करता येतील
  हे सहज समजेल असे आपण गृहीत धरतो.
\end{editorial}
\end{document}
